\section{Introduction}
\label{sec:intro}

Permissioned Byzantine Fault Tolerant (BFT) consensus underpins a growing class
of e-governance systems: electronic voting, land and business registries,
digital identity. Deploying one at national scale requires answering a concrete
question before any node is procured: \emph{how many validators should the
network have?}

The validator count~$\n$ acts in two opposite directions. Throughput argues for
small~$\n$: every validator executes every transaction and PBFT-family
protocols exchange $O(\n^{2})$ messages per block. Safety argues for
large~$\n$: detecting a coordinated faulty minority is a statistical estimation
problem whose concentration bounds tighten exponentially in~$\n$. The two
regimes have so far been studied separately; their joint treatment is the
subject of this paper.

Two obstacles stand in the way. First, \emph{scaling exponents are not
reproducible}: published power-law fits $\TPS(\n)\propto\n^{\alpha}$ disagree in
magnitude and sometimes in sign, and our own two earlier studies are an
instance --- extrapolated to sixty validators they differ by roughly an order of
magnitude. One cause is mundane. Client-side offered concurrency is rarely held
fixed while~$\n$ is varied, since a larger network is usually driven by a larger
load generator, so the fitted exponent mixes a client-parallelism effect with
the consensus-overhead effect. Second, \emph{large testbeds are unavailable}:
calibrating a scaling law is assumed to require a testbed of the target size,
and running many validators on one host confounds consensus overhead with
contention for the host CPU --- precisely the quantity one is trying to isolate.

\paragraph{Contributions.}
The first result is central; the other three exist because it is not useful on
its own.

\begin{enumerate*}
\item \textbf{An identifiability theorem} (Section~\ref{sec:model}). If
throughput obeys $\TPS(\n,\cc)=\Tzero\,\gfun(\cc)\,\n^{-\beta}$ with $\beta>0$
and the campaign varies concurrency together with the network along a path
$\log\cc=\lambda\log\n+\mathrm{const}$, then the single-factor least-squares fit
converges in probability to $\ahat=\gamma\lambda-\beta$
(Theorem~\ref{thm:identifiability}). A positive single-factor exponent is thus
evidence about~$\lambda$, not about the protocol, and exponents from different
testbeds are incomparable unless their concurrency paths are --- which is rarely
reported.

\item \textbf{A measurement protocol that makes it usable}
(Section~\ref{sec:model}). Setting $\lambda=0$ requires a factorial design, and
on one host $\n$ validators sharing $Q$ cores receive a per-validator budget
that itself depends on~$\n$. We fix a CPU quota~$\q$ per validator irrespective
of~$\n$ and report $\TPSn=\TPS/\q$; the resulting invariance
(Statement~\ref{st:qinvariance}) is tested, not assumed.

\item \textbf{Prediction intervals validated by coverage}
(Section~\ref{sec:sizing}), replacing the average percentage error that this
literature usually reports and that cannot be checked.

\item \textbf{Sizing as a chance-constrained program}
(Section~\ref{sec:sizing}). The throughput and detection constraints move in
opposite directions, so the feasible set is a closed interval $[\nmin,\nmax]$
(Theorem~\ref{thm:window}) located in $O(\log\n)$ evaluations; an empty interval
certifies that no validator-set size satisfies the requirement.
\end{enumerate*}

We claim nothing new for the detection bound itself: it is Hoeffding's
inequality with a design-effect correction, and its role is to supply the second
constraint with parameters we measure rather than postulate.

\paragraph{Scope.}
We model throughput and detection reliability only; cryptographic design,
endpoint security and coercion resistance are outside the scope. Absolute
throughput values are specific to the constrained environment, and every claim
concerns exponents and the sizing decisions derived from them. All measurements
come from one implementation on one host, which is a real limitation;
Section~\ref{sec:case} states the cross-platform prediction that
Theorem~\ref{thm:identifiability} makes and that a second implementation can
falsify.
